\question[3]
\begin{parts}
\part Eine sortierte Liste hat 1000, 1 Million bzw. 1 Milliarde Elemente. Wie viele Schleifendurchläufe
braucht die binäre Suche höchstens? Begründe kurz.
\part Wie viele Elemente muss die lineare Suche bei einer Milliarde Elementen höchstens ansehen?
\part Was liefert die binäre Suche aus dem Unterricht für die Liste \texttt{[5, 1, 4, 2, 3]} und \texttt{x = 1}?
Warum?
\end{parts}
\begin{solutionbox}{5cm}
\begin{enumerate}[label=\alph*)]
\item 10, 20 bzw. 30. Jeder Durchlauf halbiert den Suchbereich, und $1000 \approx 2^{10}$,
1 Million $\approx 2^{20}$, 1 Milliarde $\approx 2^{30}$.
\item 1 Milliarde, wenn \texttt{x} am Ende steht oder nicht in der Liste ist.
\item \texttt{-1}, obwohl 1 in der Liste steht: \texttt{mid = 2}, \texttt{a[2] = 4 > 1}, also wird links
weitergesucht; \texttt{mid = 0}, \texttt{a[0] = 5 > 1}, dann ist \texttt{R = -1}. Die Liste ist nicht sortiert,
die binäre Suche funktioniert nur für sortierte Listen.
\end{enumerate}
\end{solutionbox}
